Through systematic searches during the past decades, astronomers have found a
large number of millisecond pulsars (spin period $< 10$\,ms). The majority of
these pulsars are found in binaries, with nearly circular orbits.

For a pulsar in a binary orbit, the measured pulsar spin period ($P$) and the
measured line-of-sight acceleration ($a$) both vary systematically due to
orbital motion. For circular orbits, this variation can be described
mathematically in terms of orbital phase $\phi$ ($0 \le \phi \le 2\pi$) as,
\[ P(\phi) = P_0 + P_{\mathrm{t}}\cos\phi \qquad \mbox{where }
   P_{\mathrm{t}} = \frac{2\pi P_0 r}{c P_{\mathrm{B}}}, \]
\[ a(\phi) = -a_{\mathrm{t}}\sin\phi \qquad \mbox{where }
   a_{\mathrm{t}} = \frac{4\pi^2 r}{P_{\mathrm{B}}^2}, \]
where $P_{\mathrm{B}}$ is the orbital period of the binary, $P_0$ is the
intrinsic spin period of the pulsar and $r$ is the radius of the orbit.

The following table gives one such set of measurements of $P$ and $a$ at
different heliocentric epochs, $T$, expressed in truncated Modified Julian
Days (tMJD), i.e.\ number of days since $\mathrm{MJD} = 2\,440\,000$.

\begin{center}
\begin{tabular}{cccc}
No. & $T$ & $P$ & $a$ \\
 & (tMJD) & ($\mu$s) & (m\,s$^{-2}$) \\
1 & 5740.654 & 7587.8889 & $-0.92 \pm 0.08$ \\
2 & 5740.703 & 7587.8334 & $-0.24 \pm 0.08$ \\
3 & 5746.100 & 7588.4100 & $-1.68 \pm 0.04$ \\
4 & 5746.675 & 7588.5810 & $+1.67 \pm 0.06$ \\
5 & 5981.811 & 7587.8836 & $+0.72 \pm 0.06$ \\
6 & 5983.932 & 7587.8552 & $-0.44 \pm 0.08$ \\
7 & 6005.893 & 7589.1029 & $+0.52 \pm 0.08$ \\
8 & 6040.857 & 7589.1350 & $+0.00 \pm 0.04$ \\
9 & 6335.904 & 7589.1358 & $+0.00 \pm 0.02$ \\
\end{tabular}
\end{center}

By plotting $a(\phi)$ as a function of $P(\phi)$, we can obtain a parametric
curve. As evident from the relations above, this curve in the
period-acceleration plane is an ellipse.

In this problem, we estimate the intrinsic spin period, $P_0$, the orbital
period, $P_{\mathrm{B}}$, and the orbital radius, $r$, by an analysis of this
data set, assuming a circular orbit.

\begin{description}
\item[(D1.1)] Plot the data, including error bars, in the
  period-acceleration plane (mark your graph as ``D1.1'').
\item[(D1.2)] Draw an ellipse that appears to be a best fit to the data (on
  the same graph ``D1.1'').
\item[(D1.3)] From the plot, estimate $P_0$, $P_{\mathrm{t}}$ and
  $a_{\mathrm{t}}$, including error margins.
\item[(D1.4)] Write expressions for $P_{\mathrm{B}}$ and $r$ in terms of
  $P_0$, $P_{\mathrm{t}}$, $a_{\mathrm{t}}$.
\item[(D1.5)] Calculate approximate value of $P_{\mathrm{B}}$ and $r$ based
  on your estimations made in (D1.3), including error margins.
\item[(D1.6)] Calculate orbital phase, $\phi$, corresponding to the epochs
  of the following five observations in the above table: data rows 1, 4, 6,
  8, 9.
\item[(D1.7)] Refine the estimate of the orbital period, $P_{\mathrm{B}}$,
  using the results in part (D1.6) in the following way:
  \begin{description}
  \item[(D1.7a)] First determine the initial epoch, $T_0$, which corresponds
    to the nearest epoch of zero orbital phase before the first observation.
  \item[(D1.7b)] The expected time, $T_{\mathrm{calc}}$, of the estimated
    orbital phase angle of each observation is given by,
    \[ T_{\mathrm{calc}} = T_0
       + \left(n + \frac{\phi}{360^\circ}\right) P_{\mathrm{B}}, \]
    where $n$ is the number of full cycles of orbital phases that may have
    elapsed between $T_0$ and $T$ (or $T_{\mathrm{calc}}$). Estimate $n$ and
    $T_{\mathrm{calc}}$ for each of the five observations in part (D1.6).
    Note down the difference $T_{\mathrm{O-C}}$ between observed $T$ and
    $T_{\mathrm{calc}}$. Enter these calculations in the table given in the
    Summary Answersheet.
  \item[(D1.7c)] Plot $T_{\mathrm{O-C}}$ against $n$ (mark your graph as
    ``D1.7'').
  \item[(D1.7d)] Determine the refined values of the initial epoch,
    $T_{0,\mathrm{r}}$, and the orbital period, $P_{\mathrm{B,r}}$.
  \end{description}
\end{description}
